% Research derivation. This file is an input fragment, not a standalone paper.
% The plain-moment theorem below has a compiled Lean counterpart.
% The final assembly of the zero-density statement is a written proof.
\section{The unmarked plain moment and density for a fixed character dictionary}
\label{sec:actual-unmarked-density}

This section closes the zero-slot plain-moment input for a fixed finite
character dictionary. It does not identify the varying rows in the proposed
physical moment with such a dictionary; that separate identification must
preserve every Euler hole and the stated constants.

Let $\mathcal O=\mathbb Z[\omega]$ be the Eisenstein integers and put
$m_0=2(1-\omega)$. A free row means a nonzero element $u$ whose principal ideal
has every prime valuation less than six. Fix a nonempty finite set $\Lambda$ (the empty case is immediate) and
characters $\eta_\lambda$, $\lambda\in\Lambda$. For each $u$ choose the
\emph{natural} row character $\chi_{u,\lambda}$, so that
\[
 c_{\chi_{u,\lambda}}(I)
 =c_{\eta_\lambda}(I)\,\mathfrak s_{m_0^6u}(I),
 \qquad
 \mathrm N\mathfrak q_{u,\lambda}
 \le K_\lambda\,\mathrm N(u),
 \quad K_\lambda=\mathrm N\mathfrak q_{\eta_\lambda}
                   \mathrm N(m_0)\mathrm N(72).
\]
Here $\mathfrak s$ is the literal extended sextic ideal character, including
its zero values at ramified primes. The conductor estimate refers to the
natural representative obtained from the conductor-bounded row constructor.
The abstract \path{RowData.character} is chosen by a different existence
theorem and does not carry this bound. Its coefficients agree with the
natural representative, which is all the moment estimates require.
Whole-character conjugation preserves the modulus and reflects the zero
height. It can be handled through this symmetry; it must not be confused
with conjugating only the fixed factor $\eta_\lambda$.

\subsection{The zero-slot plain input}

For $n\le2$, $\sigma\in[0,1]$, $t\in\mathbb R$, write
\[
 P_n(\chi;D,\sigma,t)
 =D^{-1/2}\sum_{I\ne0}c_\chi(I)
   (\log(\mathrm NI/D))^n W(\mathrm NI/D)
   (\mathrm NI/D)^{-\sigma+it},
\]
where $W$ is the fixed positive annular detector profile, supported in
$[1/4,9/4]$. Its logarithmic derivatives occurring here are those of the
existing detector construction.

\paragraph{Plain-moment proposition.}
Assume the contradiction branch $\beta\ge51/100$. For each $\epsilon>0$
there are $C>0$, $J\in\mathbb N$, and $U_0$, depending on the fixed
dictionary and $\epsilon$, such that for every $U\ge U_0$, every height
$H\ge0$, every unmarked detector fiber whose right length satisfies $m\le3/4$,
and whose finite row set $\mathcal R$ satisfies $\mathrm N(u)\le U$ and
the displayed coefficient dictionary, one has
\[
 \sum_{u\in\mathcal R}
 |P_j(\chi_{u,\lambda};U^m,\sigma,t)
  P_k(\chi_{u,\lambda};U^m,\sigma,t)|^2
 \le C(1+H)^J U^{\max(1,2m)+\epsilon},
 \tag{P}
\]
for $j+k\le2$, $\sigma\in[0,1]$ and $|t|\le H$.
The threshold precedes the height quantifier. No lower bound on
$\mathrm N(u)$ and no prime-slot supply is assumed. In the formal theorem
$\mathcal R$ is the row set of an unmarked detector fiber; the proof of its
plain estimate only uses its row bound, coefficient dictionary, and $m$.

\paragraph{Proof.}
Fix a label and choose, for example,
$Q=\mathfrak q_\eta(72)$. Then $Q\ne0$, $Q\ne\mathcal O$,
$Q\subseteq(72)$, and $Q\subseteq\mathfrak q_\eta$. These are the exact
fixed-modulus gates of the incumbent analytic induction. For
$d=\min(\epsilon/4,1/4)>0$, its terminal certificate with
$\kappa=\max(3/4,2\beta-1)$ supplies the zero-slot bound at error $d$.
The zero slot profile is a legitimate input to that certificate; no
improvement of the physical moment is used. The fixed character modulus
satisfies $\mathrm N\mathfrak q_\eta\le U^d$ for all sufficiently large $U$.
The natural initial state has width
\[
 1+d+\max(0,2m-1)=\max(1,2m)+d\le2.
\]
Consequently the incumbent zero-slot estimate, followed by the uniform
Schwartz-profile bounds, proves (P) on rows satisfying
\path{initialKeep}, with exponent $\max(1,2m)+2d$.
The profile control contributes $C(1+H)^J$ and does not change the scale
threshold. The legacy proof interface is reached by adding one auxiliary
slot of width $7/37$ with identically zero profile and always selecting the
empty product. Every legacy structure field, including its numerical
supply field, is then proved; no prime sum or supply property is imposed
on the original family.

The omitted free rows form a fixed finite set $E_\eta$. Indeed, the
fixed-inducing-row theorem puts all their prime factors in the finite
support of $Q$, and free-row valuations are at most five. There are only
finitely many resulting ideals and six possible units per principal ideal.
This is a fixed finite set, not the whole interval of small rows.

The exceptional contribution has the required size without cancellation.
The detector Schwartz profiles have a uniform absolute bound $A$ for all
$n\le2$, $\sigma\in[0,1]$, and all frequencies. Since
$\#\{I:0<\mathrm NI\le x\}\le128x$, coefficient norms are at most one,
and the profile vanishes above $9/4$, the normalized finite sum gives
\[
 |P_n(\chi;D,\sigma,t)|\le288A\sqrt D
\]
for every character, including principal characters. Therefore
\[
 \sum_{u\in\mathcal R\cap E_\eta}|P_jP_k|^2
 \le |E_\eta|(288A)^4D^2
 \le |E_\eta|(288A)^4U^{\max(1,2m)+\epsilon}.
\]
This proves (P) on all rows. Taking maxima of finitely many thresholds and
degrees, and a sum of the finitely many constants, proves the simultaneous
statement for $\Lambda$. The compiled theorem
\path{fixed_unmarked_plain} proves this argument for one label; finite
label closure is an ordinary finite maximum argument.

\subsection{From actual witnesses to a height-uniform density bound}

Let $N(Y,a,T)$ count free rows $u$, $\mathrm N(u)\le Y$, for which at least
one dictionary character has a zero with real part at least $a$ and
imaginary part of absolute value at most $T$. In this paragraph
$17/20\le a\le7/8$. Use the incumbent bound $\beta\le7/8$. If $\beta<a$, the
nonexceptional contribution is zero; the fixed exceptional set is handled
by its constant cardinality. The lower hypothesis $\beta\ge51/100$
in (P) is automatic whenever the nonexceptional contribution is nonempty
and hence is not an extra theorem assumption on the desired zero-free result.

\paragraph{Density proposition.}
For every $\varepsilon>0$ there are $C,J,Y_0$, fixed before $a$ and $T$,
such that for all $Y\ge Y_0$ and $T\ge0$,
\[
 N(Y,a,T)\le C(1+T)^JY^{\rho(2a-1)+\varepsilon},
 \qquad
 \rho(\delta)=1-\frac{10\delta}{3(5-2\delta)}.
 \tag{D}
\]
The constant depends on the fixed dictionary. The same assertion for a
scale-dependent dictionary does not follow from this proposition.

\paragraph{Proof.}
Set $K=\max(1,\max_\lambda K_\lambda)$ and $U=KY$. Then both the row norm
and every natural character modulus are at most $U$. This is a fixed
multiplicative padding; no power of $Y$ is lost. Remove the finite union
$E=\bigcup_\lambda E_{\eta_\lambda}$. On the remaining rows every natural
character is nonprincipal, by the incumbent nonprincipal gate. The finite
union contributes $O(1)$ to (D), so there is no need to construct detector
witnesses for its members.

Choose small detector and moment errors before the scale parameter. The
actual inverse moment for each fixed \path{momentData} family and (P)
now give both moment inputs with one bound $C(1+H)^{J_0}$. Both the constant
and degree are independent of the detector zero bin, the dyadic indices,
and $U$. The actual detector lengths satisfy
$m\le1/2+75\epsilon_{\rm det}\le23/40$, so (P) applies.
Choose $\tau_{\rm det}>0$ \emph{after} $J_0$, sufficiently small for the supported
witness construction and for
$\tau_{\rm det}(J_0+1)$ to fit within the allotted final error.

First suppose $T\le U^{\tau_{\rm det}}$. For each row, apply the existing finite
buffered-bin lemma to its finite character family with height step
$U^{\tau_{\rm det}}>2$ and mesh $e>0$. It gives an integer $1\le i\le I$,
$I=\lceil49/(100e)\rceil$, and a real grid value $b=51/100+ek$ such that
\[
 b\le M(3iU^{\tau_{\rm det}})<b+e,
 \qquad M(3(i+1)U^{\tau_{\rm det}})<b+2e.
\]
Here $M$ is the actual maximum of zero real parts in the specified box,
with the usual $51/100$ fallback. Because the original row has a zero
with real part at least $a$ at height at most $T$, monotonicity gives
$b>a-e$. The incumbent upper bound gives $b\le7/8$. There are only
$O_e(1)$ possibilities for $(i,k)$, so partition the rows by these values.
For each part set $\delta_b=2b-1$ and fix its balanced cutoff
$t_*=5/(5-2\delta_b)$ before constructing witnesses.

The actual supported-witness theorem, with base scale $Z=U$ and $d=1$,
produces witnesses in each nonempty part. Its hypotheses are now literal:
nonprincipality, modulus at most $U$, the two buffered maxima, and the
natural-row coefficient identity. It uses frequency allowance
$U^{\tau_{\rm det}/2}$, and the witness frequency is bounded by
\[
 H_i=2\pi U^{\tau_{\rm det}/2}+3iU^{\tau_{\rm det}}\ll_e U^{\tau_{\rm det}}.
\]
Choose the auxiliary small parameters in the stated detector budget
$12e(22+2)+8\kappa_{\rm aux}+2\eta_{\rm aux}
 \le\epsilon_{\rm det}/2$ and
$4\tau_{\rm det}<\eta_{\rm aux}$. They can all be chosen positively, with
$e<1/1000$, and then $\tau_{\rm det}$ decreased further as necessary. No height
variable occurs in the resulting eventual threshold.

Partition once more by the witness label and its two dyadic indices, all
constructed using this fixed $t_*$. The no-slot inverse and plain estimates yield the
already compiled balanced count
\[
 |\mathcal R_{\rm fiber}|
 \ll (1+H_i)^{J_0+1}
 U^{\rho(\delta_b)+156\epsilon_{\rm det}+\epsilon_m}.
\]
The fixed label count and the two dyadic counts cost an arbitrary positive
power of $U$. Since $b>a-e$ and $\rho$ is decreasing, differentiation on
the relevant compact interval gives
$\rho(2b-1)\le\rho(2a-1)+3e$.
Choose $\epsilon_{\rm det},\epsilon_m,e$ and the dyadic padding so their
sum with $\tau_{\rm det}(J_0+1)$ is less than $\varepsilon$. Summing the finite
partitions proves $N(Y,a,T)\ll U^{\rho(2a-1)+\varepsilon}$ in this case.
The replacement $U=KY$ only changes the fixed constant.

If $T>U^{\tau_{\rm det}}$, the elementary lattice count is enough:
$N(Y,a,T)\ll U$. Since $\rho(2a-1)>0$ and $U\ge1$, for any integer
$J\ge1/\tau_{\rm det}$,
\[
 U\le T^J\le(1+T)^JU^{\rho(2a-1)+\varepsilon}.
\]
This proves (D) for arbitrary $T$, with a common threshold fixed before
that height. Bounded $Y$ are covered by enlarging the constant using the
finite row ball, uniformly in $T$ and $a$.

When this result is used at an external physical height $T=Z^{\tau_{\rm out}}$,
choose $\tau_{\rm out}$ only after the final degree $J$ in (D) has been
fixed, so that $J\tau_{\rm out}$ fits the physical error allowance. The
external exponent $\tau_{\rm out}$ is distinct from the internal detector
exponent $\tau_{\rm det}$ chosen in this proof. Reusing a single exponent
for both purposes would make the parameter selection circular.

The exponent satisfies $\rho(\delta)\ge2/7>1/6$ for
$\delta\le3/4$. Thus the weighted sixth-power raw-row lifting proved in
\path{HeightUniformDensity.lean} applies to (D). This completes the
analytic unmarked density input for a fixed dictionary. Applying it to the
joint physical row still requires the exact signed coefficient identity,
a fixed-dictionary reduction for every varying label, and the stated
conductor and mask accounting.

\subsection{Formal verification scope}

The new module \path{PlainUnmarkedAdapter.lean} contains nine audited
objects: the natural-row coefficient dictionary, the legacy embedding, the retained estimate, the actual zero-slot
input, the retained eventual field, finiteness of exceptional free rows,
the absolute detector-polynomial bound, the finite-row absolute bound, and
the complete eventual plain field. All nine printed axiom sets contain
only \path{propext}, \path{Classical.choice}, and \path{Quot.sound}.
Compilation reuses development dependency objects; this is not a clean
source replay or an independent checker certificate. The assembly of (D)
above is a written mathematical proof, not an additional Lean theorem in
this module.
